\documentclass[11pt,letterpaper]{article}
% Reusable style. Build each .tex from its own directory using scripts/build.sh.
\usepackage[T1]{fontenc}
\usepackage[utf8]{inputenc}
\usepackage{lmodern}
\usepackage[margin=0.86in,headheight=15pt]{geometry}
\usepackage{amsmath,amssymb,amsthm,mathtools}
\usepackage{booktabs,longtable,tabularx,array,multicol}
\usepackage[table]{xcolor}
\usepackage{microtype}
\usepackage{enumitem}
\usepackage{fancyhdr}
\usepackage{titlesec}
\usepackage{hyperref}
\usepackage{listings}
\usepackage[most]{tcolorbox}
\usepackage{needspace}
\definecolor{aeNavy}{HTML}{172D43}
\definecolor{aeBlue}{HTML}{176B90}
\definecolor{aeTeal}{HTML}{148477}
\definecolor{aePale}{HTML}{F2F7F9}
\definecolor{aeGray}{HTML}{5C6771}
\hypersetup{colorlinks=true,linkcolor=aeBlue,urlcolor=aeBlue,pdfauthor={Algorithm Engineering Documentation Suite},pdfsubject={Independent educational materials based on the supplied TAOCP volume topics}}
\titleformat{\section}{\Large\bfseries\color{aeNavy}}{\thesection}{0.6em}{}
\titleformat{\subsection}{\large\bfseries\color{aeBlue}}{\thesubsection}{0.6em}{}
\titleformat{\subsubsection}{\normalsize\bfseries\color{aeNavy}}{\thesubsubsection}{0.6em}{}
\setlist[itemize]{leftmargin=*,topsep=3pt,itemsep=1pt}
\setlist[enumerate]{leftmargin=*,topsep=3pt,itemsep=3pt}
\setlength{\parindent}{0pt}\setlength{\parskip}{5pt}
\pagestyle{fancy}\fancyhf{}\fancyhead[L]{\footnotesize\color{aeGray}ALGORITHM ENGINEERING}\fancyhead[R]{\footnotesize\color{aeGray}DOCUMENTATION SUITE}\fancyfoot[L]{\footnotesize\color{aeGray}Independent study and engineering practice}\fancyfoot[R]{\footnotesize\thepage}
\renewcommand{\headrulewidth}{0.35pt}
\newtcolorbox{engineeringnote}[1][]{breakable,colback=aePale,colframe=aeBlue,boxrule=0.5pt,arc=1mm,fonttitle=\bfseries,title={Engineering note},#1}
\newtcolorbox{evidencebox}[1][]{breakable,colback=white,colframe=aeTeal,boxrule=0.6pt,arc=1mm,fonttitle=\bfseries,title={Evidence and deliverables},#1}
\lstdefinestyle{algo}{basicstyle=\ttfamily\footnotesize,breaklines=true,frame=single,rulecolor=\color{aeGray!40},backgroundcolor=\color{aePale},showstringspaces=false,tabsize=2,columns=fullflexible,keepspaces=true,keywordstyle=\color{aeBlue}\bfseries,commentstyle=\color{aeTeal}}
\lstset{style=algo}
\newcommand{\DocTitle}[3]{%
\hypersetup{pageanchor=false}\begin{titlepage}\vspace*{1.2cm}
{\color{aeTeal}\rule{\linewidth}{2.3pt}}\par\vspace{1.15cm}
{\fontsize{24}{30}\selectfont\bfseries\color{aeNavy}\raggedright #1\par}\vspace{0.45cm}
{\Large\color{aeBlue}#2\par}\vspace{1.3cm}
{\large\bfseries #3\par}\vspace{0.35cm}
{\color{aeGray}Edition 1.0 \quad | \quad 29 October 2026\par}\vfill
{\color{aeGray}\small Original educational material aligned to the chapter and section topics in the five supplied volumes of Donald E. Knuth's \emph{The Art of Computer Programming}. This document is not an authorized reproduction or substitute for the original books.}\par\vspace{0.4cm}
{\color{aeTeal}\rule{\linewidth}{0.75pt}}
\end{titlepage}\hypersetup{pageanchor=true}\pagenumbering{roman}}
\newcommand{\SourceNote}[1]{\par\smallskip{\footnotesize\color{aeGray}\textbf{Source locator:} #1}\par\smallskip}
\newcommand{\Outcome}[1]{\par\textbf{Competency target.} #1\par}
\newcommand{\Checkpoint}[1]{\par\textbf{Checkpoint.} #1\par}
% Intentionally simple printable writing lines: call \AnswerLines{number}.
\newcount\alcount
\newcommand{\AnswerLines}[1]{\par\alcount=0\loop\ifnum\alcount<#1\par\noindent\rule{\linewidth}{0.3pt}\vspace{0.33cm}\advance\alcount by1\repeat\par}
\newcommand{\FormField}[2][0.62\linewidth]{\noindent\textbf{#2}\hspace{0.65em}\rule{#1}{0.4pt}\par\vspace{0.18cm}}
\newcommand{\ReviewTable}{\begin{tabularx}{\linewidth}{@{}p{0.23\linewidth}X p{0.2\linewidth}@{}}\toprule Criterion & Recorded evidence & Rating (0--4)\\\midrule Correctness & & \\[1em] Proof quality & & \\[1em] Complexity & & \\[1em] Testing & & \\[1em] Reproducibility & & \\[1em] Communication & & \\[1em]\bottomrule\end{tabularx}}

\begin{document}
\DocTitle{Volume 2 Study Guide}{Chapters 3--4: randomness and arithmetic}{Chapter- and section-aligned study guide}
\tableofcontents
\clearpage
\pagenumbering{arabic}
\section{Study protocol}

For each subsection, read the specified part of the original volume, keep a mathematical notebook, produce an original implementation or design sketch, and evaluate the acceptance gate. The investigations below are new exercises, not copies of the book's problems or solutions. Annotate any starred or mathematically advanced source sections as optional on the first pass; revisit them after the prerequisite mathematics is in place.
\begin{evidencebox}Recommended artifacts: study notes, problem specification, derivation, runnable code or model, tests, performance evidence and reflective comparison to original assumptions.\end{evidencebox}

\section{\S 3.1--3.2: Uniform pseudorandom generation}
\SourceNote{Knuth, Volume 2, \S 3.1--3.2.}
\Outcome{Understand deterministic state updates, modulus, recurrence period, seeding and portability.}
\subsection{Independent investigation}
Implement a configurable linear congruential generator and enumerate its full cycle for a small modulus.
\subsection{Analysis and verification gate}
Report observed period; do not claim cryptographic suitability.
\subsection{Study questions}
\begin{enumerate}
\item Which representations, input assumptions and cost metrics control this technique's applicability?
\item Which invariant or combinatorial identity establishes correctness, and which input is its strongest counterexample candidate?
\item What result would cause the implementation to fail review despite passing a small collection of examples?
\end{enumerate}
\Checkpoint{Preserve a short explanation, the relevant test evidence, and one unresolved question before proceeding.}

\section{\S 3.3: Empirical, theoretical and spectral tests}
\SourceNote{Knuth, Volume 2, \S 3.3.}
\Outcome{Separate marginal uniformity from serial dependence; understand type-I error and sample size.}
\subsection{Independent investigation}
Collect digit frequencies, adjacent pairs and a chi-squared statistic on fixed-seed samples.
\subsection{Analysis and verification gate}
Interpret significance cautiously, including false rejection risk and low-count limitations.
\subsection{Study questions}
\begin{enumerate}
\item Which representations, input assumptions and cost metrics control this technique's applicability?
\item Which invariant or combinatorial identity establishes correctness, and which input is its strongest counterexample candidate?
\item What result would cause the implementation to fail review despite passing a small collection of examples?
\end{enumerate}
\Checkpoint{Preserve a short explanation, the relevant test evidence, and one unresolved question before proceeding.}

\section{\S 3.4: Transforms, shuffling and sampling}
\SourceNote{Knuth, Volume 2, \S 3.4.}
\Outcome{Derive uniform integer mapping, unbiased rejection sampling and Fisher--Yates shuffling.}
\subsection{Independent investigation}
Compare Fisher--Yates with a deliberately biased swap scheme on small permutations.
\subsection{Analysis and verification gate}
Use a frequency table and justify expected probabilities, not visual impression.
\subsection{Study questions}
\begin{enumerate}
\item Which representations, input assumptions and cost metrics control this technique's applicability?
\item Which invariant or combinatorial identity establishes correctness, and which input is its strongest counterexample candidate?
\item What result would cause the implementation to fail review despite passing a small collection of examples?
\end{enumerate}
\Checkpoint{Preserve a short explanation, the relevant test evidence, and one unresolved question before proceeding.}

\section{\S 3.5--3.6: Meaning of randomness}
\SourceNote{Knuth, Volume 2, \S 3.5--3.6.}
\Outcome{Distinguish statistical test suites, algorithmic randomness and cryptographic unpredictability.}
\subsection{Independent investigation}
Draft a threat-model-aware selection guideline for simulation seeds versus session tokens.
\subsection{Analysis and verification gate}
Reject any inference from passing frequency tests to adversarial unpredictability.
\subsection{Study questions}
\begin{enumerate}
\item Which representations, input assumptions and cost metrics control this technique's applicability?
\item Which invariant or combinatorial identity establishes correctness, and which input is its strongest counterexample candidate?
\item What result would cause the implementation to fail review despite passing a small collection of examples?
\end{enumerate}
\Checkpoint{Preserve a short explanation, the relevant test evidence, and one unresolved question before proceeding.}

\section{\S 4.1--4.2: Positional and floating-point arithmetic}
\SourceNote{Knuth, Volume 2, \S 4.1--4.2.}
\Outcome{Analyze bases, rounding, cancellation, overflow and representable spacing.}
\subsection{Independent investigation}
Construct a summation workload with severe cancellation and compare different summation orders.
\subsection{Analysis and verification gate}
Record absolute/relative error against a higher precision reference where possible.
\subsection{Study questions}
\begin{enumerate}
\item Which representations, input assumptions and cost metrics control this technique's applicability?
\item Which invariant or combinatorial identity establishes correctness, and which input is its strongest counterexample candidate?
\item What result would cause the implementation to fail review despite passing a small collection of examples?
\end{enumerate}
\Checkpoint{Preserve a short explanation, the relevant test evidence, and one unresolved question before proceeding.}

\section{\S 4.3--4.4: Multiple precision and radix conversion}
\SourceNote{Knuth, Volume 2, \S 4.3--4.4.}
\Outcome{Model multiprecision cost in digits and validate base conversion algorithms.}
\subsection{Independent investigation}
Convert random integers between bases 2, 10 and 16; round-trip against language built-ins.
\subsection{Analysis and verification gate}
Test zero, large magnitudes and leading-zero conventions.
\subsection{Study questions}
\begin{enumerate}
\item Which representations, input assumptions and cost metrics control this technique's applicability?
\item Which invariant or combinatorial identity establishes correctness, and which input is its strongest counterexample candidate?
\item What result would cause the implementation to fail review despite passing a small collection of examples?
\end{enumerate}
\Checkpoint{Preserve a short explanation, the relevant test evidence, and one unresolved question before proceeding.}

\section{\S 4.5: Rationals, GCD, factorization}
\SourceNote{Knuth, Volume 2, \S 4.5.}
\Outcome{Establish Euclidean invariants and define normalized rational representation.}
\subsection{Independent investigation}
Implement Euclid and binary GCD with a differential test harness.
\subsection{Analysis and verification gate}
Explain termination and compare operation counts, not only wall time.
\subsection{Study questions}
\begin{enumerate}
\item Which representations, input assumptions and cost metrics control this technique's applicability?
\item Which invariant or combinatorial identity establishes correctness, and which input is its strongest counterexample candidate?
\item What result would cause the implementation to fail review despite passing a small collection of examples?
\end{enumerate}
\Checkpoint{Preserve a short explanation, the relevant test evidence, and one unresolved question before proceeding.}

\section{\S 4.6: Polynomial arithmetic}
\SourceNote{Knuth, Volume 2, \S 4.6.}
\Outcome{Use Horner evaluation, polynomial division assumptions and exponentiation-by-squaring.}
\subsection{Independent investigation}
Implement Horner and naive evaluation using a bounded integer domain.
\subsection{Analysis and verification gate}
Demonstrate equal results, different multiplication counts and overflow caveats.
\subsection{Study questions}
\begin{enumerate}
\item Which representations, input assumptions and cost metrics control this technique's applicability?
\item Which invariant or combinatorial identity establishes correctness, and which input is its strongest counterexample candidate?
\item What result would cause the implementation to fail review despite passing a small collection of examples?
\end{enumerate}
\Checkpoint{Preserve a short explanation, the relevant test evidence, and one unresolved question before proceeding.}

\section{\S 4.7: Power series}
\SourceNote{Knuth, Volume 2, \S 4.7.}
\Outcome{Understand coefficient-wise formal series manipulations and truncation.}
\subsection{Independent investigation}
Compute truncated series products and compare low-order coefficients to a trusted small calculation.
\subsection{Analysis and verification gate}
Make truncation degree and coefficient-domain assumptions explicit.
\subsection{Study questions}
\begin{enumerate}
\item Which representations, input assumptions and cost metrics control this technique's applicability?
\item Which invariant or combinatorial identity establishes correctness, and which input is its strongest counterexample candidate?
\item What result would cause the implementation to fail review despite passing a small collection of examples?
\end{enumerate}
\Checkpoint{Preserve a short explanation, the relevant test evidence, and one unresolved question before proceeding.}

\section{Volume synthesis and mastery review}
\begin{enumerate}
\item Draw a dependency graph connecting at least four major ideas from this volume.
\item Select two alternative implementations for one workload and compare their preconditions, proven cost, memory use, measured performance and failure cases.
\item Present a five-minute technical defense of the most surprising experimental result. A claim that differs from the mathematical model requires a measured or documented explanation.
\end{enumerate}
\subsection{Completion checklist}
\begin{itemize}\item All module notes include a source locator and an independently stated learning outcome.
\item Every implementation has reproducible tests and edge cases.
\item Every complexity claim identifies its cost model and underlying assumptions.
\item A peer can independently reproduce one chosen experiment using its recorded configuration.
\end{itemize}
\end{document}
